\section{Seed-Wise OA Results}
\label{sec:appendix_seedwise}

For transparency, Table~\ref{tab:seedwise_results} reports the per-seed OA values of the baseline, the full method, and the key ablation variants that support the stability discussion in the main text. This appendix table is especially useful for checking whether a variant improves the mean through consistent gains or through a small number of favorable runs.

\begin{table}[t]
  \centering
  \caption{Per-seed OA values on the fixed split manifests used by all compared methods.}
  \label{tab:seedwise_results}
  \begin{tabular}{l l c c c}
    \toprule
    Dataset & Variant & Seed 3407 & Seed 3408 & Seed 3409 \\
    \midrule
    AID-low20 & ConvNeXt-Tiny & 0.9232 & 0.9058 & 0.9248 \\
    AID-low20 & DINOv2-B (Adapter, FT) & 0.9365 & 0.9423 & 0.9452 \\
    AID-low20 & NoCenter & 0.9411 & 0.9395 & 0.9244 \\
    AID-low20 & NoAdapter & 0.9326 & 0.9438 & 0.9485 \\
    NWPU-low20 & ConvNeXt-Tiny & 0.8860 & 0.8726 & 0.8575 \\
    NWPU-low20 & DINOv2-B (Adapter, FT) & 0.8863 & 0.8835 & 0.8671 \\
    NWPU-low20 & NoCenter & 0.8768 & 0.8643 & 0.8755 \\
    NWPU-low20 & NoAdapter & 0.8946 & 0.8753 & 0.8680 \\
    \bottomrule
  \end{tabular}
\end{table}


\section{Reproducibility Notes}
\label{sec:appendix_reproducibility}

All compared methods use seed-controlled manifests generated by class-wise shuffling under the same split rule: 20 training images and 10 validation images per class, with all remaining images used for testing. For a fixed seed, the same manifest is reused by every compared method, so the ConvNeXt baseline and all DINOv2 variants are evaluated on identical data partitions. The final release package will include the split-generation logic, model definitions, training loop, multi-seed aggregation utilities, and the scripts used to rebuild the manuscript tables.

The final DINOv2 setting uses the configuration corresponding to the 20-shot protocol, partial finetuning of the last transformer block, adapter bottleneck dimension 256, and feature-center loss weight 0.02. The final ConvNeXt baseline uses the same seed protocol and early-stopping logic but a lighter backbone and larger batch size. No seed-specific hyperparameter changes are introduced after the protocol is fixed. The accompanying reproducibility package also retains the manuscript build script and the submission-readiness checker used to regenerate the final C\&G submission PDF.
